\documentclass[12pt,a4paper]{article}

\input{/home/tabaka/Documents/GitHub/Defs/defs.tex}

\setcounter{zadaniaRozdzial}{1}

\begin{document}
	
	\ZbiorZadanieTwoXTwo{Oblicz wartości funkcji trygonometrycznych kąta $AOP$, jeśli $O$ to początek układu współrzędnych, $A$ leży na dodatniej części osi $OX$, a punkt $P$ jest taki, że:}
	{$P=(3,4)$}
	{$P=(6,8)$}
	{$P=(1,2)$}
	{$P=(\sqrt{3},1)$}
	
	\ZbiorZadanieTwoXTwo{Oblicz $\cos^2\alpha$, $\cos\alpha$, $\tg\alpha$ dla kąta $\alpha\in(90^\circ,180^\circ)$ spełniającego warunek:}
	{$\sin\alpha=\frac{3}{5}$}
	{$\sin\alpha=\frac{2}{3}$}
	{$\ctg\alpha=-\frac{1}{2}$}
	{$\ctg\alpha=-\sqrt{3}$}
	
	\ZbiorZadanieTwoXTwo{Wyznacz pozostałe funkcje trygonometryczne kąta $\alpha$, jeśli:}
	{$\sin\alpha=\frac{12}{13}$, $\alpha\in(0^\circ,90^\circ)$}
	{$\tg\alpha=3$, $\alpha\in(180^\circ,270^\circ)$}
	{$\cos\alpha=\frac{\sqrt{5}}{3}$, $\alpha\in(270^\circ,360^\circ)$}
	{$\ctg\alpha=-2$, $\alpha\in(90^\circ,180^\circ)$}
	
	\ZbiorZadanieTwoXThree{Oblicz podane wartości trygonometryczne:}
	{$\sin 120^\circ$}
	{$\cos 150^\circ$}
	{$\tg 135^\circ$}
	{$\sin 210^\circ$}
	{$\cos 300^\circ$}
	{$\ctg 315^\circ$}
	
	\WzorElastyczny{
		\begin{center}
			\begin{tabular}{p{5cm} p{5cm} p{5cm}}
				$\sin(90^\circ-\alpha)=\cos\alpha$
				&
				$\cos(90^\circ-\alpha)=\sin\alpha$
				&
				\\[6pt]
				$\sin(90^\circ+\alpha)=\cos\alpha$
				&
				$\cos(90^\circ+\alpha)=-\sin\alpha$
				&
				\\[6pt]
				$\sin(180^\circ-\alpha)=\sin\alpha$
				&
				$\cos(180^\circ-\alpha)=-\cos\alpha$
				&
				$\tg(180^\circ-\alpha)=-\tg\alpha$
				\\[6pt]
				$\sin(180^\circ+\alpha)=-\sin\alpha$
				&
				$\cos(180^\circ+\alpha)=-\cos\alpha$
				&
				$\tg(180^\circ+\alpha)=\tg\alpha$
			\end{tabular}
		\end{center}
	}
	\newpage
	\ZbiorZadanieTwoXThree{Oblicz podane wartości trygonometryczne:}
	{$\sin 240^\circ$}
	{$\cos 330^\circ$}
	{$\tg 225^\circ$}
	{$\ctg 150^\circ$}
	{$\sin 315^\circ$}
	{$\cos 390^\circ$}
	
	\ZbiorZadanieTwoXThree{Oblicz podane wartości trygonometryczne:}
	{$\sin 420^\circ$}
	{$\cos 750^\circ$}
	{$\tg 495^\circ$}
	{$\sin 840^\circ$}
	{$\cos 1020^\circ$}
	{$\ctg 765^\circ$}
	
	\ZbiorZadanieTwoXThree{Oblicz:}
	{$\sin 120^\circ+\cos 300^\circ$}
	{$\cos 150^\circ-\sin 210^\circ$}
	{$\tg 225^\circ+\cos 330^\circ$}
	{$\sin 240^\circ-\cos 120^\circ$}
	{$\cos 315^\circ+\sin 150^\circ$}
	{$\tg 135^\circ-\ctg 300^\circ$}
	
	\ZbiorZadanieTwoXThree{Oblicz:}
	{$\sin 120^\circ\cdot\cos 300^\circ$}
	{$\cos 150^\circ\cdot\sin 210^\circ$}
	{$\tg 225^\circ\cdot\cos 330^\circ$}
	{$\sin 240^\circ\cdot\cos 120^\circ$}
	{$\cos 315^\circ\cdot\sin 150^\circ$}
	{$\tg 135^\circ\cdot\ctg 300^\circ$}
	
	\ZbiorZadanieTwoXThree{Oblicz:}
	{$\frac{\sin 120^\circ- \cos 150^\circ}{3\tg 150^\circ}$}
	{$\frac{\tg 120^\circ-\cos 30^\circ}{\sin 120^\circ}$}
	{$\frac{\tg 150^\circ-\tg 120^\circ}{\cos 120^\circ}$}
	{$\frac{\tg 135^\circ+\cos 45^\circ}{\sin 45^\circ \cdot \cos135^\circ}$}
	{$\frac{\cos 135^\circ + \cos 150^\circ}{\sin 120^\circ + \sin 135^\circ}$}
	{$\frac{\tg 120^\circ-\cos30^\circ}{\cos 150^\circ+\cos 60^\circ}$}
	
	\ZbiorZadanieTwoXThree{Korzystając z tablic wartości trygonometrycznych, oblicz przybliżoną wartość:}{$\sin 105^\circ$}{$\sin 165^\circ$}{$\cos 130^\circ$}{$\cos 175^\circ$}{$\tg 142^\circ$}{$\tg 103^\circ$}
	
\end{document}